\documentclass[../main.tex]{subfiles}

\begin{document}
\subsection{Menghitung Suhu Ruangan}
\subsubsection{Algoritma}
\begin{enumerate}
    \item []
    \begin{enumerate}
        \item {Mengambil masukan berupa suhu awal.}
        \item {Menampilakan suhu awal dalam Celsius, Fahrenheit dan Kelvin}
        \item {Menampilkan 15 menit berlalu, menurunkan suhu sebanyak 25\%. dari suhu awal.}
        \item {Menampilkan suhu saat ini dalam Fahrenheit dan selisih suhu saat ini dengan suhu awal dalam Kelvin.}
        \item {Menampilkan 45 menit berlalu, menurunkan suhu sebanyak 65\% dari suhu awal}
        \item {Menampilkan suhu saat ini dalam Kelvin dan selisih suhu saat ini dengan suhu awal dalam Celsius.}
        \item {Mematikan pendingin dan menunjukan 25 menit berlalu, menaikan suhu sebanyak 65\% dari suhu terakhir dan menampilkan suhu saat ini.}
        \item {Menanpilkan perbandingan suhu awal dengan saat ini, perbandingan suhu awal dengan suhu terdingin dan perbandingan suhu awal dan suhu setelah 15 menit.}
    \end{enumerate}
\end{enumerate}

\subsubsection{{\Pseudocode}}
\inputcode[n]{res/m1p1.txt}

\subsubsection{{\Flowchart}}
\inputfigure[n]{res/m1f1.png}{0.43}{{\Flowchart} Program Menghitung Suhu Ruangan}{fig:m1:f1}

\subsubsection{{\Source} {\sCode}}
\inputcode[n]{res/m1c1.cpp}

\subsubsection{Hasil Program}
\inputfigure[n]{res/m1i1.png}{0.25}{Tampilan Program Menghitung Suhu Ruangan}{fig:m1:i1}

\begin{indentpar}
    \hspace{\parindent}Berdasarkan \Cref{fig:m1:i1} di atas, didapatkan keluaran dari program Menghitung Suhu Ruangan. Di baris teratas program meminta masukan berupa suhu kepada pengguna, kemudian mengeluarkannya di paragraf selanjutnya. Pada paragraf tersebut ditampilkan suhu masukan di dalam tiga unit yaitu Celsius, Fahrenheit dan Kelvin.

    Paragraf seterusnya menunjukan perubahan waktu dan dampaknya pada suhu awal. Seperti yang dapat dilihat pada paragraf ketiga semakin banyak waktu berlalu semakin menurun suhu ruangan. Di paragraf setelah itu ditunjukan suhu awal ruangan dan suhu saat ini setelah 45 menit berlalu dari awal jalannya program.

    Baris selanjutnya menunjukan keadaan pendingin yang telah dimatikan. Paragraf setelah itupun menunjukan suhu ruangan setelah dimatikan yang telah naik dibandingkan suhu 25 menit sebelumnya saat pendingin baru dimatikan.

    Paragraf terakhir menunjukan perbandingan suhu ruangan awal dengan suhu ruangan saat ini, suhu ruangan terdingin dan suhu ruangan 15 menit setelah pendinginan yang dimana disana dapat kita lihat skala perbedaan suhu ruangannya.
\end{indentpar}

\end{document}
